%%--------------------------------------------------------------------
%1--------------------------------------------------------------------
\begin{glyph}{Read (読)}{read}\label{glyph:read}
Read.
\end{glyph}
\gindex{Read}{読}\strindex{14}{読}

\begin{comps}
\comprtwocone & 言 + 売 \\\hline
\comprthreecone & Say/Speaking/Speech + Sell (\ref{glyph:sell}) \\\hline
\end{comps}

\begin{remglyph}
\hooglig{Reading} material, they \remcom{say}, \remcom{sells} a lot.\\\hline
\end{remglyph}

\begin{prons}
\pronsrtwocone & \textbf{ドク (DOKU)} & \textbf{よ (yo)}\\\hline
\pronsrthreecone & トク (TOKU)  & --- \\\hline
\end{prons}

\begin{remprons}
\hooglig{Reading} about the \ucomon{talking} \comkun{yomp} moving \comon{dockward}.\\\hline
\end{remprons}

%{\middel}
% &  ()\\\hline
\begin{somme}
読む  & \textit{to read} & よ{\middel}む \mbox{(yo{\middel}mu)} \\\hline
読み切れ  & read + cut = \textit{to finish reading} & よ{\middel}み き{\middel}る \mbox{(yo{\middel}mi ki{\middel}re)} \\\hline
読者  & read + inidividual = \textit{reader}  & ドク{\middel}シャ \mbox{(DOKU{\middel}SHA)} \\\hline
\notyet{読み物}  & read + thing = \textit{reading matter}   & よみ{\middel} もの \mbox{(yo{\middel}mi mono)} \\\hline
\notyet{多読家} & many + read + house = \textit{well-read person} & タ{\middel}ドク{\middel}シャ \mbox{(TA{\middel}DOKU{\middel}KA)} \\\hline
\end{somme}

\begin{decomp}{6}
此の&古い&本&を&読んでいます&か。\\\hline
この&ふるい&ホン&を&よんでいます&か\\\hline
this&old&book&を&are reading&か\\\hline
\multicolumn{6}{|r|}{\textit{Are you reading this old book?}}\\\hline
\end{decomp}

\end{multicols}
%%--------------------------------------------------------------------
%2--------------------------------------------------------------------
\begin{multicols}{2}
\begin{glyph}{Road (道)}{road}\label{glyph:road}
Road.  Way.\\\hline
This interesting 漢字 can signify both a \textit{physical road} and a spiritual \textit{path}.  In this latter sense, it is often translated as \textit{the way}.
\end{glyph}
\gindex{Road}{道}\strindex{12}{道}

\begin{comps}
\comprtwocone & 首 + 辵辶⻌⻍  \\\hline
\comprthreecone & Neck/Head + Road/Walk \\\hline
\end{comps}

\begin{remglyph}
\hooglig{Road} travelling with my \remcom{neck}-support.\\\hline
\end{remglyph}

\begin{prons}
\pronsrtwocone & \textbf{ドウ (D\={O})} & \textbf{みち (michi)}\\\hline
\pronsrthreecone & トウ (T\={O})  & ---  \\\hline
\end{prons}

\begin{remprons}
\hooglig{Roads} in \comkun{Michigan} are so \ucomon{torturous} that you'd rather open the \comon{door}, step out, and start walking.\\\hline
\end{remprons}

%{\middel}
% &  ()\\\hline
\begin{somme}
道 & \textit{road; way}  & みち \mbox{(michi)}  \\\hline
小道 & small + road = \textit{path}  & こ{\middel}みち \mbox{(ko{\middel}michi)}  \\\hline
水道 & water + door = \textit{water supply system}  & スイ{\middel}ドウ \mbox{(SUI{\middel}D\={O})}  \\\hline
下水道 & lower + water + road = \textit{sewer system}  & ゲ{\middel}スイ{\middel}ドウ \mbox{(GE{\middel}SUI{\middel}D\={O})} \\\hline
歩道 & walk + road = \textit{sidewalk}  & ホ{\middel}ドウ \mbox{(HO{\middel}D\={O})}  \\\hline
国道 & country + road = \textit{national highway}  & コク{\middel}ドウ \mbox{(KOKU{\middel}T\={O})}  \\\hline
\notyet{神道} & God + road = \textit{Japanese religion}  & シン{\middel}トウ \mbox{(SHIN{\middel}T\={O})} \\\hline
\end{somme}

\begin{decomp}{5}
右&の&道&を&行け。\\\hline
みち&の&みち&を&いけ\\\hline
right&の&road&を&to go\\\hline
\multicolumn{5}{|r|}{\textit{Take the right road.}}\\\hline
\end{decomp}

\end{multicols}
%%--------------------------------------------------------------------
%3--------------------------------------------------------------------
\begin{multicols}{2}
\begin{glyph}{Heaven (天)}{heaven}\label{glyph:heaven}
Heaven.  The Heavens.\\\hline
This character can indicate both the spiritual sense of \textit{heaven} as well as the more literal notions of \textit{sky} and \textit{air}.
\end{glyph}
\gindex{Heaven}{天}\strindex{4}{天}

\begin{comps}
\comprtwocone & 大 + 一 \\\hline
\comprthreecone & Large + One \\\hline
\end{comps}

\begin{remglyph}
\hooglig{Heaven} is \remcom{one} \remcom{large} place for the saints.\\\hline
\end{remglyph}

\begin{prons}
\pronsrtwocone & \textbf{テン (TEN)} & ---\\\hline
\pronsrthreecone & ---  & あめ、あま (ame, ama) \\\hline
\end{prons}

\begin{remprons}
\hooglig{Heaven} is not \comon{tense} because there is no need for \ucomkun{amends} or \ucomkun{amalgamations}.\\\hline
\end{remprons}

%{\middel}
% &  ()\\\hline
\begin{somme}
\rye{3}{Note how the 音読み for \textit{country} is voiced in the first example.  In the fourth compound, the 音読み of 王 changes from オウ to ノウ in order to make the pronunciation smoother.  This type of phonetic change is extremely rare.}
天国  & heaven + country = \textit{heaven}  & テン{\middel}コク \mbox{(TEN{\middel}GOKU)} \\\hline
天下  & heaven + lower = \textit{the whole land}  & テン{\middel}カ \mbox{(TEN{\middel}KA)} \\\hline
天体  & heaven + body = \textit{heavenly body}  & テン{\middel}タイ \mbox{(TEN{\middel}TAI)}  \\\hline
天王星  & heaven + king + star = \textit{Uranus}  & テン{\middel}ノウ{\middel}セイ \mbox{(TEN{\middel}N\={O}{\middel}SEI)}  \\\hline
\notyet{天気}  & heaven + spirit = \textit{weather}  & テン{\middel}キ \mbox{(TEN{\middel}KI)}  \\\hline
\end{somme}

\begin{decomp}{8}
明日&の&天気&は&雨&だ&そう&です。\\\hline
あした&の&テンキ&は&あめ&だ&そう&です\\\hline
tomorrow&の&weather&は&rain&だ&seems&be / is\\\hline
\multicolumn{8}{|r|}{\textit{It seems the weather will be rainy tomorrow.}}\\\hline
\end{decomp}

\end{multicols}
%%--------------------------------------------------------------------
%4--------------------------------------------------------------------
\begin{multicols}{2}
\begin{glyph}{Part (分)}{part}\label{glyph:part}
This important character suggests the ideas of \textit{dividing} and \textit{portion}.  The sense of \textit{division} is always present in some way.\\\hline
The first compound below can be though of as \textit{parting} the unknown to arrive at understanding.
\end{glyph}
\gindex{Part}{分}\strindex{4}{分}

\begin{comps}
\comprtwocone & 八 + 刀 \\\hline
\comprthreecone & Eight + Sword \\\hline
\end{comps}

\begin{remglyph}
\hooglig{Part} the \remcom{volcano} with this \remcom{sword}.\\\hline
\end{remglyph}

\begin{prons}
\pronsrtwocone & \textbf{ブン (BUN)} & \textbf{わ (wa)}\\\hline
\pronsrthreecone & フン、ブ (FUN, BU) & ---  \\\hline
\rye{3}{フン relates to minutes (of time or arc), as in the example 五分.
ブ deals with percentages and rates; 五分 in this sense would be pronounced ゴブ, and mean either five or fifty percent.\\\hline
Note also that 大分, one of the many Japanese words expressing the idea of \textit{much} or \textit{very}, can be pronounced ダイブン or ダイブ.}
\end{prons}

\begin{remprons}
The \hooglig{part} where I \comkun{wander} to the \comon{boon} such that I \ucomon{foon}.\\\hline
\end{remprons}

%{\middel}
% &  ()\\\hline
\begin{somme}
分かる (intr.) & \textit{to understand}  & わ{\middel}かる \mbox{(wa{\middel}karu)}  \\\hline
分ける (tr.) & \textit{to part; to divide}  & わ{\middel}ける \mbox{(wa{\middel}keru)}  \\\hline
自分 & self + part = \textit{oneself}  & ジ{\middel}ブン \mbox{(JI{\middel}BUN)}  \\\hline
半分 & half + part = \textit{half}  & ハン{\middel}ブン \mbox{(HAN{\middel}BUN)}  \\\hline
春分 & spring + part = \textit{vernal equinox}  & シュン{\middel}ブン \mbox{(SHUN{\middel}BUN)}  \\\hline
秋分 & autumn + part = \textit{autumnal equinox}  & シュウ{\middel}ブン \mbox{(SH\={U}{\middel}BUN)}  \\\hline
五分 & five + part = \textit{five minutes}  & ゴフ{\middel}ン \mbox{(GO{\middel}FUN)}  \\\hline
五分 & five + part = \textit{five/fifty percent}  & ゴ{\middel}ブ \mbox{(GO{\middel}BU)}  \\\hline
大分 & large + part = \textit{much/very}  & ダイ{\middel}ブン、ダイ{\middel}ブ \mbox{(DAI{\middel}BU(N))}  \\\hline
\end{somme}

\begin{decomp}{4}
日本語&が&分かります&か。\\\hline
ニホンゴ&が&わかります&か\\\hline
Japanese&が&understand&か\\\hline
\multicolumn{4}{|r|}{\textit{Do you understand Japanese.}}\\\hline
\end{decomp}

\end{multicols}
\pagebreak
%%--------------------------------------------------------------------
%5--------------------------------------------------------------------

\begin{multicols}{2}
\begin{glyph}{Thing (物)}{thing}\label{glyph:thing}
Thing.  Object.\\\hline
Note that the writing order for \textit{cow} differs when it is used as a component; slanting the final stroke upward makes for a smoother transition to the next part of the character.
\end{glyph}
\gindex{Thing}{物}\strindex{8}{物}

\begin{comps}
\comprtwocone & 牛 + 勿 \\\hline
\comprthreecone & Cow + (\ldots) Not \\\hline
\end{comps}

\begin{remglyph}
That \hooglig{thing} can \remcom{not} be a \remcom{cow}.\\\hline
\end{remglyph}

\begin{prons}
\pronsrtwocone & \textbf{ブツ (BUTSU)} & \textbf{もの (mono)}\\\hline
\pronsrthreecone & モツ (MOTSU)  & --- \\\hline
\end{prons}

\begin{remprons}
The \hooglig{thing} about this \comkun{monologue} is \ucomon{motsurella} \comon{boots}.\\\hline
\end{remprons}

%{\middel}
% &  ()\\\hline
\begin{somme}
物 & \textit{thing}  & もの \mbox{(mono)} \\\hline
買い物 & buy + thing = \textit{shopping}  & か{\middel}い　もの \mbox{(ka{\middel}i mono)} \\\hline
読み物 & read + thing = \textit{reading matter}  & よ{\middel}み　もの \mbox{(yo{\middel}mi mono)} \\\hline
人物 & person + thing = \textit{one's character}  & ジン{\middel}ブツ \mbox{(JIN{\middel}BUTSU)} \\\hline
見物 & see + thing = \textit{sightseeing}  & ケン{\middel}ブツス \mbox{(KEN{\middel}BUTSU)} \\\hline
物体 & thing + body = \textit{an object}  & ブッ{\middel}タイ \mbox{(BUT{\middel}TAI)} \\\hline
\end{somme}

\begin{decomp}{6}
本田さん&は&どんな&人物&です&か。\\\hline
ホンださん&は&どんな&ジンブツ&です&か\\\hline
Honda-san&は&what kind of&character&be / is&か\\\hline
\multicolumn{6}{|r|}{\textit{What kind of person is Honda-san?}}\\\hline
\end{decomp}

\end{multicols}

%%--------------------------------------------------------------------
%6--------------------------------------------------------------------
\begin{multicols}{2}
\begin{glyph}{Axe (斤)}{axe}\label{glyph:axe}
Axe.\\\hline
Though rarely encountered on its own, this character serves as a component in several important 漢字.
\end{glyph}
\gindex{Axe}{斤}\strindex{4}{斤}

\begin{comps}
\comprtwocone & 斤 \\\hline
\comprthreecone & Axe/Kind of Weight Measure \\\hline
\end{comps}

\begin{remglyph}
Part of the 漢字 radicals (\#{69}).\\\hline
\end{remglyph}

\begin{prons}
\pronsrtwocone & \textbf{キン (KIN)} & ---\\\hline
\pronsrthreecone & --- & --- \\\hline
\end{prons}

\begin{remprons}
\hooglig{Axe} the \comon{kindred}.\\\hline
\end{remprons}

%{\middel}
% &  ()\\\hline
\begin{somme}
斤 & \textit{axe}  & キン \mbox{(KIN)}  \\\hline
\end{somme}

\begin{decomp}{4}
斤&が&読んでます&か。\\\hline
キン&が&よんでます&か\\\hline
axe&が&read&か\\\hline
\multicolumn{4}{|r|}{\textit{Can you read ``斤''?}}\\\hline
\end{decomp}

\end{multicols}

\pagebreak
%%--------------------------------------------------------------------
%7--------------------------------------------------------------------

\begin{multicols}{2}
\begin{glyph}{New (新)}{new}\label{glyph:new}
New.
\end{glyph}
\gindex{New}{新}\strindex{13}{新}

\begin{comps}
\comprtwocone & 立 + 木 + 斤 \\\hline
\comprthreecone & Stand + Tree + Axe\\\hline
\end{comps}

\begin{remglyph}
\hooglig{New} \remcom{trees} cannot \remcom{stand} \remcom{axes}.\\\hline
\end{remglyph}

\begin{prons}
\pronsrtwocone & \textbf{シン (SHIN)} & \textbf{あたら (atara)}\\\hline
\pronsrthreecone & ---  & あら、にい (ara, n\={i}) \\\hline
\end{prons}

\begin{remprons}
\hooglig{New} \ucomkun{arabesques} will make you \comon{sheen} with \ucomkun{neat} \comkun{ataraxia}.\\\hline
{\warn}\textit{Arabesque} is a ballet position.  {\warn}\textit{Ataraxia} is a state of serene calmness.\\\hline
\end{remprons}

%{\middel}
% &  ()\\\hline
\begin{somme}
新しい & \textit{new}  & あたら{\middel}しい \mbox{(atara{\middel}shii)} \\\hline
一新 & one + new = \textit{complete renewal}  & イッ{\middel}シン \mbox{(IS{\middel}SHIN)} \\\hline
新月 & new + moon = \textit{new moon}  & シン{\middel}ゲツ \mbox{(SHIN{\middel}GETSU)} \\\hline
新星 & new + star = \textit{nova}  & シン{\middel}セイ \mbox{(SHIN{\middel}SEI)}  \\\hline
\notyet{最新} & most + new = \textit{newest}  & サイ{\middel}シン \mbox{(SAI{\middel}SHIN)} \\\hline
\end{somme}

\begin{decomp}{4}
新しい&日本刀&を&買いました。\\\hline
あたらしい&ニホントウ&を&かいました\\\hline
new&Japanese sword&を&bought\\\hline
\multicolumn{4}{|r|}{\textit{(I) bought a new Japanese sword.}}\\\hline
\end{decomp}

\end{multicols}

%%--------------------------------------------------------------------
%8--------------------------------------------------------------------
\begin{multicols}{2}
\begin{glyph}{Flower (花)}{flower}\label{glyph:flower}
Flower.
\end{glyph}
\gindex{Flower}{花}\strindex{7}{花}

\begin{comps}
\comprtwocone & 艸 + 化 \\\hline
\comprthreecone & Grass + Change (\ref{glyph:change}) \\\hline
\end{comps}

\begin{remglyph}
\hooglig{Flowers} are \remcom{grass} \remcom{changelings}.\\\hline
\end{remglyph}

\begin{prons}
\pronsrtwocone & \textbf{カ (KA)} & \textbf{はな (hana)}\\\hline
\pronsrthreecone & --- & --- \\\hline
\end{prons}

\begin{remprons}
\hooglig{Flowers} that \comon{cuss} are \comkun{hung-up}.\\\hline
\end{remprons}

%{\middel}
% &  ()\\\hline
\begin{somme}
花 & \textit{flower}  & はな \mbox{(hana)} \\\hline
生花 & life + flower = \textit{flower arrangement}  & い{\middel}け　ばな \mbox{(i{\middel}ke bana)} \\\hline
花見 & flower + see = \textit{cherry blossom viewing}  & はな{\middel}み \mbox{(hana{\middel}mi)} \\\hline
花火 & flower + fire = \textit{fireworks}  & はな{\middel}び \mbox{(hana{\middel}bi)} \\\hline
火花 & fire + flower = \textit{sparks}  & ひ{\middel}ばな \mbox{(hi{\middel}bana)} \\\hline
開花 & open + flower = \textit{bloom}  & カイ{\middel}カ \mbox{(KAI{\middel}KA)}  \\\hline
\end{somme}

\begin{decomp}{5}
お&花見&に&行きません&か。\\\hline
お&はなみ&に&いきません&か\\\hline
お&flower viewing&に&to go&か\\\hline
\multicolumn{5}{|r|}{\textit{Why don't we go and see the cherry blossoms?}}\\\hline
\end{decomp}

\end{multicols}

\pagebreak
%%--------------------------------------------------------------------
%9--------------------------------------------------------------------

\begin{multicols}{2}
\begin{glyph}{Name (名)}{name}\label{glyph:name}
Name.\\\hline
By extension, this 漢字 also includes the ideas of \textit{reputation} and \textit{fame}.
\end{glyph}
\gindex{Name}{名}\strindex{6}{名}

\begin{comps}
\comprtwocone & 夕 + 口\\\hline
\comprthreecone & Evening + Mouth\\\hline
\end{comps}

\begin{remglyph}
\hooglig{Names} of \remcom{evening} \remcom{vampires}.\\\hline
\end{remglyph}

\begin{prons}
\pronsrtwocone & \textbf{メイ (MEI)} & \textbf{な (na)}\\\hline
\pronsrthreecone & ミョウ (MY\={O}) & --- \\\hline
\end{prons}

\begin{remprons}
\hooglig{Naming} the \ucomon{Myou} dolls is a \comkun{numbing} \comon{mayem}.\\\hline
\end{remprons}

%{\middel}
% &  ()\\\hline
\begin{somme}
名高い  & name + tall = \textit{renowned}  & な　だ{\middel}かい \mbox{(na da{\middel}kai)} \\\hline
有名  & have + name = \textit{famous}  & ユウ{\middel}メイ \mbox{(Y\={U}{\middel}MEI)} \\\hline
名山  & name + mountain = \textit{famous mountain}  & メイ{\middel}ザン \mbox{(MEI{\middel}ZAN)} \\\hline
名物  & name + thing = \textit{famous regional product}  & メイ{\middel}ブツ \mbox{(MEI{\middel}BUTSU)}  \\\hline
\notyet{名前}  & name + before = \textit{name}  & な{\middel}まえ \mbox{(na{\middel}mae)}  \\\hline
\end{somme}

\begin{decomp}{5}
彼女&は&有名&歌手&だ。\\\hline
かのジョ&は&ユウメイ&カシュ&だ\\\hline
she&は&famous&singer&be / is\\\hline
\multicolumn{5}{|r|}{\textit{She is a noted singer.}}\\\hline
\end{decomp}

\end{multicols}

%%--------------------------------------------------------------------
%10-------------------------------------------------------------------
\begin{multicols}{2}
\begin{glyph}{Mother (母)}{mother}\label{glyph:mother}
Mother.
\end{glyph}
\gindex{Mother}{母}\strindex{5}{母}

\begin{comps}
\comprtwocone & 母 \\\hline
\comprthreecone & Mother \\\hline
\end{comps}

\begin{remglyph}
Part of the 漢字 radicals (\#{80}).\\\hline
\end{remglyph}

\begin{prons}
\pronsrtwocone & \textbf{ボ (BO)} & \textbf{はは (haha)}\\\hline
\pronsrthreecone & --- & --- \\\hline
\end{prons}

\begin{remprons}
\hooglig{Mother} \comon{bottles} up all her \comkun{laughter}.\\\hline
\end{remprons}

%{\middel}
% &  ()\\\hline
\begin{somme}
母 & \textit{mother (responsible occupation)}  &  はは \mbox{(haha)}\\\hline
母親 & mother + parent = \textit{mother (formal)}  &  はは{\middel}おや \mbox{(haha{\middel}oya)}\\\hline
母国語 & mother + country + words = \textit{mother tongue}  & ボ{\middel}コク{\middel}ゴ \mbox{(BO{\middel}KOKU{\middel}GO)} \\\hline
母の日 & mother + sun (day) = \textit{Mother's day}  &  はは{\middel}の{\middel}ひ \mbox{(haha{\middel}no{\middel}hi)} \\\hline
父母 & father + mother = \textit{father and mother}  &  フ{\middel}ボ \mbox{(FU{\middel}BO)} \\\hline
\end{somme}

\begin{irrr}
お母さん  & \textit{mother}  &  おかあさん (ok\={a}san) \\\hline
伯母さん{\notcov}  & related + mother = \textit{aunt}  & おばさん (obasan) \\\hline
\end{irrr}

\end{multicols}
\begin{decomp}{7}
新しい&先生&の&母国語&は&英語&です。\\\hline
あたらしい&センセイ&の&ボコクゴ&は&エイゴ&です\\\hline
new&teacher&の&mother tongue&は&English&be / is\\\hline
\multicolumn{7}{|r|}{\textit{The new teacher's mother tongue is English.}}\\\hline
\end{decomp}

\pagebreak
%%--------------------------------------------------------------------
%11-------------------------------------------------------------------

\begin{multicols}{2}
\begin{glyph}{Private (私)}{private}\label{glyph:private}
Private.\\\hline
This character is also used for the first person singular, \textit{I}.
\end{glyph}
\gindex{Private}{私}\strindex{7}{私}

\begin{comps}
\comprtwocone & 禾 + 厶 \\\hline
\comprthreecone & Two-branch tree/Grain + Private \\\hline
\end{comps}

\begin{remglyph}
\hooglig{\remcom{Private}} \remcom{grain}.\\\hline
\end{remglyph}

\begin{prons}
\pronsrtwocone & \textbf{シ (SHI)} & \textbf{わたし (watashi)}\\\hline
\pronsrthreecone & --- & わたくし (watakushi) \\\hline
\end{prons}

\begin{remprons}
I \hooglig{privately} rejoiced--``\comkun{what a shield}, \ucomkun{what a
cushion}''.\\\hline
\end{remprons}

%{\middel}
% &  ()\\\hline
\begin{somme}
私 & \textit{I}  & わたし \mbox{(watashi)} \\\hline
私有 & private + have = \textit{privately owned}  & シ{\middel}ユウ \mbox{(SHI{\middel}Y\={U})} \\\hline
私道 & private + road = \textit{private road}  & シ{\middel}ドウ \mbox{(SHI{\middel}D\={O})} \\\hline
私物 & private + thing = \textit{private property}  & シ{\middel}ブツ \mbox{(SHI{\middel}BUTSU)} \\\hline
\notyet{私学}  & private + study = \textit{private school}  & シ{\middel}ガク \mbox{(SHI{\middel}GAKU)}  \\\hline
\end{somme}

\begin{decomp}{3}
私&を&愛してる？\\\hline
わたし&を&アイしてる\\\hline
I&を&to love\\\hline
\multicolumn{3}{|r|}{\textit{Do you love me?}}\\\hline
\end{decomp}

\end{multicols}

%%--------------------------------------------------------------------
%12-------------------------------------------------------------------
\begin{multicols}{2}
\begin{glyph}{Night (夜)}{night}\label{glyph:night}
Night.
\end{glyph}
\gindex{Night}{夜}\strindex{8}{夜}

\begin{comps}
\comprtwocone & 亠 + 人亻𠆢 + 夕 \\\hline
\comprthreecone & Top/Lid + Person + Evening\\\hline
\end{comps}

\begin{remglyph}
\hooglig{Night} \remcom{people} wear \remcom{tops} when it is \remcom{evening}.\\\hline
\end{remglyph}

\begin{prons}
\pronsrtwocone & \textbf{ヤ (YA)} & \textbf{よる、よ (yoru, yo)}\\\hline
\rye{3}{よる is only used when this 漢字 appears alone.\\\hline
よ and ヤ are distributed evenly throughout compounds.}
\pronsrthreecone & --- & --- \\\hline
\end{prons}

\begin{remprons}
A \hooglig{night}-time \comkun{yomp} to the \comon{yacht} with \comkun{Yoruichi}.\\\hline
{\warn}\textit{Yoruichi} 「夜一」 is a character in the famous \textit{Bleach} \mbox{アニメ}.\\\hline
\end{remprons}

%{\middel}
% &  ()\\\hline
\begin{somme}
夜  & \textit{night}  & よる \mbox{(yoru)} \\\hline
夜中  & night + middle = \textit{middle of the night}  & よ{\middel}なか \mbox{(yo{\middel}naka)} \\\hline
夜明け  & night + bright = \textit{dawn}  & よ　あ{\middel}け \mbox{(yo a{\middel}ke)} \\\hline
月夜  & moon + night = \textit{moonlit night}  & つき{\middel}よ \mbox{(tsuki{\middel}yo)} \\\hline
夜間  & night + interval = \textit{night time}  & ヤ{\middel}カン \mbox{(YA{\middel}KAN)} \\\hline
\end{somme}

\begin{decomp}{5}
夜明け&に&家&を&出ました。\\\hline
よあけ&に&いえ&を&でました\\\hline
dawn&に&house&を&exit\\\hline
\multicolumn{5}{|r|}{\textit{I left the house at dawn.}}\\\hline
\end{decomp}

\end{multicols}

\pagebreak
%%--------------------------------------------------------------------
%13-------------------------------------------------------------------

\begin{multicols}{2}
\begin{glyph}{Before (前)}{before}\label{glyph:before}
Before.  Front.\\\hline
It is worth comparing the compounds below to those for 先 (Entry \ref{glyph:precede}, page \pageref{glyph:precede}) to get a sense of the subtle differences between these two characters.
\end{glyph}
\gindex{Before}{前}\strindex{9}{前}

\begin{comps}
\comprtwocone & 艸艹 + 月 + 刀 \\\hline
\comprthreecone & Grass + Moon + Sword\\\hline
\end{comps}

\begin{remglyph}
\hooglig{Before} the \remcom{moon} appears I cut the \remcom{grass} with my \remcom{sword}.\\\hline
\end{remglyph}

\begin{prons}
\pronsrtwocone & \textbf{ゼン (ZEN)} & \textbf{まえ (mae)}\\\hline
\pronsrthreecone & --- & --- \\\hline
\end{prons}

\begin{remprons}
\hooglig{Before} the \comkun{maelstrom} there was \comon{zen}.\\\hline
\end{remprons}

%{\middel}
% &  ()\\\hline
\begin{somme}
前  & \textit{before}  & まえ \mbox{(mae)} \\\hline
名前  & name + before = \textit{name}  & な{\middel}まえ \mbox{(na{\middel}mae)} \\\hline
手前  & hand + before = \textit{this side of}  & て{\middel}まえ \mbox{(te{\middel}mae)}  \\\hline
前半  & before + half = \textit{the first half}  & ゼン{\middel}ハン \mbox{(ZEN{\middel}HAN)} \\\hline
前者  & before + individual = \textit{the former}  & ゼン{\middel}シャ \mbox{(ZEN{\middel}SHA)} \\\hline
\notyet{午前}  & noon + before = \textit{A.M.; Morning}  & ゴゼン \mbox{(GO{\middel}ZEN)} \\\hline
\end{somme}

\end{multicols}
\begin{decomp}{8}
門&の&前&に&美しい&花&が&あります。\\\hline
モン&の&まえ&に&うつくしい&はな&が&あります\\\hline
gate&の&before&に&beautiful&flower&が&are\\\hline
\multicolumn{8}{|r|}{\textit{There are beautiful flowers in front of the gate.}}\\\hline
\end{decomp}

%%--------------------------------------------------------------------
%14-------------------------------------------------------------------
\begin{multicols}{2}
\begin{glyph}{Many (多)}{many}\label{glyph:many}
Many.  Much.
\end{glyph}
\gindex{Many}{多}\strindex{6}{多}

\begin{comps}
\comprtwocone & 夕 + 夕 \\\hline
\comprthreecone & Evening + Evening\\\hline
\end{comps}

\begin{remglyph}
\hooglig{Many} sins are done during \remcom{evenings}.\\\hline
\end{remglyph}

\begin{prons}
\pronsrtwocone & \textbf{タ (TA)} & \textbf{おお (\={o})}\\\hline
\pronsrthreecone & --- & --- \\\hline
\end{prons}

\begin{remprons}
\hooglig{Many} \comkun{awful} \comon{tub} experiences.\\\hline
\end{remprons}

%{\middel}
% &  ()\\\hline
\begin{somme}
多い & \textit{many}  & おお{\middel}い \mbox{(\={o}{\middel}i)} \\\hline
多少 & many + few = \textit{more or less}  & タ{\middel}ショウ \mbox{(TA{\middel}SH\={O})} \\\hline
多分 & many + part = \textit{perhaps}  & タ{\middel}ブン \mbox{(TA{\middel}BUN)} \\\hline
多大 & many + large = \textit{great quantity}  & タ{\middel}ダイ \mbox{(TA{\middel}DAI)} \\\hline
多読家 & many + read + house = \textit{well-read person}  & タ{\middel}ドク{\middel}カ \mbox{(TA{\middel}DOKU{\middel}KA)}  \\\hline
\end{somme}

\begin{decomp}{6}
彼&は&多分&有名&に&成らない。\\\hline
かれ&は&タブン&ユウメイ&に&ならない\\\hline
he&は&perhaps&famous&に&not become\\\hline
\multicolumn{6}{|r|}{\textit{Perhaps he'll never become famous.}}\\\hline
\end{decomp}

\end{multicols}

\pagebreak
%%--------------------------------------------------------------------
%15-------------------------------------------------------------------

\begin{multicols}{2}
\begin{glyph}{Hold (持)}{hold}\label{glyph:hold}
Hold.  Possess.
\end{glyph}
\gindex{Hold}{持}\strindex{9}{持}

\begin{comps}
\comprtwocone & 手扌 + 寺 \\\hline
\comprthreecone & Hand + Temple (\ref{glyph:temple}) \\\hline
\end{comps}

\begin{remglyph}
\hooglig{Hold} the \remcom{temple} in your \remcom{hand}.\\\hline
\end{remglyph}

\begin{prons}
\pronsrtwocone & --- & \textbf{も (mo)}\\\hline
\pronsrthreecone & ジ (JI) & --- \\\hline
\end{prons}

\begin{remprons}
\hooglig{Hold} on to your \comkun{morales}, or the \ucomon{jig} is up.\\\hline
\end{remprons}

%{\middel}
% &  ()\\\hline
\begin{somme}
\rye{3}{As the fourth example below shows, this 漢字 can act like 買.}
持つ  & \textit{to hold}  & もつ \mbox{(motsu)}  \\\hline
金持ち  & gold + hold = \textit{rich person}  & かね　も{\middel}ち \mbox{(kane mo{\middel}chi)} \\\hline
持ち上げる  & hold + upper = \textit{to lift up}  & も{\middel}ち　あ{\middel}げる \mbox{(mo{\middel}chi a{\middel}geru)} \\\hline
持物  & hold + thing = \textit{one's belongings}  & もち{\middel}もの \mbox{(mochi{\middel}mono)} \\\hline
\notyet{気持ち}  & spirit + hold = \textit{feeling}  & キ　も{\middel}ち \mbox{(KI mo{\middel}chi)} \\\hline
\end{somme}

\begin{decomp}{4}
気持ち&が&悪く&なった。\\\hline
キもち&が&わるく&なった\\\hline
feeling&が&bad&to \textit{become}\\\hline
\multicolumn{4}{|r|}{\textit{A strange feeling came over me.}}\\\hline
\end{decomp}

\end{multicols}

%%--------------------------------------------------------------------
%16-------------------------------------------------------------------
\begin{multicols}{2}
\begin{glyph}{Write (書)}{write}\label{glyph:write}
Write.\\\hline
This character has to do with all matters related to writing, and can suggest such things as books, notes, documents, and calligraphy, amongst others.
\end{glyph}
\gindex{Write}{書}\strindex{10}{書}

\begin{comps}
\comprtwocone & 聿 + 日 \\\hline
\comprthreecone & Brush + Sun \\\hline
\end{comps}

\begin{remglyph}
\hooglig{Write} with this \remcom{brush} the kanji for \remcom{``Sun''}.\\\hline
\end{remglyph}

\begin{prons}
\pronsrtwocone & \textbf{ショ (SHO)} & \textbf{か (ka)}\\\hline
\pronsrthreecone & --- & --- \\\hline
\end{prons}

\begin{remprons}
\hooglig{Write} down the \comkun{customers'} \comon{shopping} experiences.\\\hline
\end{remprons}

%{\middel}
% &  ()\\\hline
\begin{somme}
\rye{3}{The third and the fifth examples below show that this 漢字 can act like 買.  Be aware that the 訓読み is always voiced when it appears in the second position, as it is in the final compound.}
書く & \textit{to write}  & か{\middel}く \mbox{(ka{\middel}ku)} \\\hline
書き入れる & write + enter = \textit{to write in}  & か{\middel}き　い{\middel}れる \mbox{(ka{\middel}ki i{\middel}reru)} \\\hline
書入れ & write + enter = \textit{an entry}  & か{\middel}き　い{\middel}れ \mbox{(ka{\middel}ki i{\middel}re)} \\\hline
書道 & write + road = \textit{calligraphy}  & ショ{\middel}ドウ \mbox{(SHO{\middel}D\={O})} \\\hline
{\diff}前書 & before + write = \textit{preface}  & まえ{\middel}がき \mbox{(mae{\middel}gaki)}  \\\hline
\end{somme}

\begin{decomp}{4}
鉛筆&で&書いて&下さい。\\\hline
エンピツ&で&かいて&ください\\\hline
pencil&で&to write&please\\\hline
\multicolumn{4}{|r|}{\textit{Please write with a pencil.}}\\\hline
\end{decomp}

\begin{decomp}{3}
手紙&書いて&ね。\\\hline
てがみ&かいて&ね\\\hline
letter&to write&\textit{request confirmation}\\\hline
\multicolumn{3}{|r|}{\textit{Write me sometime, OK?}}\\\hline
\end{decomp}

\end{multicols}

\pagebreak
%%--------------------------------------------------------------------
%17-------------------------------------------------------------------

\begin{multicols}{2}
\begin{glyph}{Sound (音)}{sound}\label{glyph:sound}
Sound.
\end{glyph}
\gindex{Sound}{音}\strindex{9}{音}

\begin{comps}
\comprtwocone & 音 \\\hline
\comprthreecone & Sound \\\hline
\end{comps}

\begin{remglyph}
Part of the 漢字 radicals (\#{180}).\\\hline
\end{remglyph}

\begin{prons}
\pronsrtwocone & \textbf{オン (ON)} & \textbf{おと (oto)}\\\hline
\pronsrthreecone & イン (IN) & ね (ne) \\\hline
\end{prons}

\begin{remprons}
\hooglig{Sounded} like an \ucomon{inkling} of a \ucomkun{necromancer} \comkun{automatically} going \comon{online}.\\\hline
\end{remprons}

%{\middel}
% &  ()\\\hline
\begin{somme}
\rye{3}{The second example is very familiar to you by now.  Note also how the meaning of the same compound changes in the example following this, when both 漢字 are read with their 音読み.}
音  & \textit{sound}  & おと \mbox{(oto)} \\\hline
音読み  & sound + read = \textit{the on-yomi}  & オン　よ{\middel}み \mbox{(ON yo{\middel}mi)} \\\hline
音読  & sound + read = \textit{reading aloud}  & オン{\middel}ドク \mbox{(ON{\middel}DOKU)} \\\hline
物音  & thing + sound = \textit{a noise}  & もの{\middel}おと \mbox{(mono{\middel}oto)} \\\hline
\notyet{足音}  & leg + sound = \textit{sound of footsteps}  & あし{\middel}おと \mbox{(ashi{\middel}oto)}  \\\hline
\end{somme}

\end{multicols}
\begin{decomp}{8}
あれ&は&水&の&音&だ&と&思います。\\\hline
あれ&は&みず&の&おと&だ&と&おもいます\\\hline
that&は&water&の&sound&だ&と&to think\\\hline
\multicolumn{8}{|r|}{\textit{I think that's the sound of water.}}\\\hline
\end{decomp}

%%--------------------------------------------------------------------
%18-------------------------------------------------------------------
\begin{multicols}{2}
\begin{glyph}{Meet (会)}{meet}\label{glyph:meet}
Meet.  Meeting.\\\hline
When found at the end of compounds, this 漢字 often expresses English words such as \textit{society}, \textit{association}, and \textit{party}.
\end{glyph}
\gindex{Meet}{会}\strindex{6}{会}

\begin{comps}
\comprtwocone & 人 + 二 + 厶 \\\hline
\comprthreecone & Person + Two + Private \\\hline
\end{comps}

\begin{remglyph}
\hooglig{Meet} \remcom{two} \remcom{people} \remcom{privately}.\\\hline
\end{remglyph}

\begin{prons}
\pronsrtwocone & \textbf{カイ (KAI)} & \textbf{あ (a)}\\\hline
\pronsrthreecone &エ (E)  &  \\\hline
\end{prons}

\begin{remprons}
\hooglig{Meet} at the \comkun{updated} \comon{kites} anchored by \ucomon{eggs}.\\\hline
\end{remprons}

%{\middel}
% &  ()\\\hline
\begin{somme}
会う  & \textit{to meet}  & あ{\middel}う \mbox{(a{\middel}u)} \\\hline
会見  & meet + see = \textit{interview}  & カイ{\middel}ケン \mbox{(KAI{\middel}KEN)} \\\hline
{\diff}大会  & large + meet = \textit{convention; tournament}  & タイ{\middel}カイ \mbox{(TAI{\middel}KAI)} \\\hline
開会  & open + meet = \textit{to open a meeting}  & カイ{\middel}カイ \mbox{(KAI{\middel}KAI)} \\\hline
入会金  & enter + meet + gold = \textit{enrollment fee}  & ニュウ{\middel}カイ{\middel}キン \mbox{(NY\={U}{\middel}KAI{\middel}KIN)}  \\\hline
\end{somme}

\begin{decomp}{6}
米国大会&は&九月&に&あります&か。\\\hline
ベイコクタイカイ&は&クガツ&に&あります&か\\\hline
U.S. tournament&は&September&に&is&か\\\hline
\multicolumn{6}{|r|}{\textit{Is the U.S. tournament in September?}}\\\hline
\end{decomp}

\end{multicols}

\pagebreak
%%--------------------------------------------------------------------
%19-------------------------------------------------------------------

\begin{multicols}{2}
\begin{glyph}{Hear (聞)}{hear}\label{glyph:hear}
Hear.\\\hline
Ask.\\\hline
Like 早 (Entry \ref{glyph:early}, page \pageref{glyph:early}), this is another 漢字 with two unrelated and distinct meanings.
\end{glyph}
\gindex{Hear}{聞}\strindex{14}{聞}

\begin{comps}
\comprtwocone & 門 + 耳\\\hline
\comprthreecone & Gate + Ear\\\hline
\end{comps}

\begin{remglyph}
\hooglig{Hearing} my \hooglig{question} starts \remcom{gate} of the \remcom{ear}.\\\hline
\end{remglyph}

\begin{prons}
\pronsrtwocone & \textbf{ブン (BUN)} & \textbf{き (ki)}\\\hline
\pronsrthreecone &モン (MON)  & --- \\\hline
\end{prons}

\begin{remprons}
\hooglig{Hear} how the \ucomon{monastery} \comkun{keeps} \comon{boons}.\\\hline
\end{remprons}

%{\middel}
% &  ()\\\hline
\begin{somme}
聞こえる (intr.)  & \textit{to be heard}  & き{\middel}こえる \mbox{(ki{\middel}koeru)}  \\\hline
聞く (tr.)  & \textit{to hear; to ask}  & き{\middel}く \mbox{(ki{\middel}ku)} \\\hline
聞き出す  & hear + exit = \textit{to find out about}  & き{\middel}き　だ{\middel}す \mbox{(ki{\middel}ki da{\middel}su)} \\\hline
聞き入れる  & hear + enter = \textit{to comply with}  & き{\middel}き　い{\middel}れる \mbox{(ki{\middel}ki i{\middel}reru)}  \\\hline
新聞  & new + hear = \textit{newspaper}  & シン{\middel}ブン \mbox{(SHIN{\middel}BUN)} \\\hline
見聞  & see + hear = \textit{information}  & ケン{\middel}ブン \mbox{(KEN{\middel}BUN)}  \\\hline
\end{somme}

\begin{decomp}{5}
夜中&に&物音&を&聞きました。\\\hline
よなか&に&ものおと&を&ききました\\\hline
middle of night&に&sound&を&heard\\\hline
\multicolumn{5}{|r|}{\textit{(She) heard a sound in the middle of the night.}}\\\hline
\end{decomp}

\end{multicols}

%%--------------------------------------------------------------------
%20-------------------------------------------------------------------
\begin{multicols}{2}
\begin{glyph}{Center (央)}{center}\label{glyph:center}
Center.\\\hline
Other than the first example (and examples derived from it), this 漢字 is rarely used.
\end{glyph}
\gindex{Center}{央}\strindex{5}{央}

\begin{comps}
\comprtwocone & 冖 + 大 \\\hline
\comprthreecone & Cover/Crown + Large\\\hline
\end{comps}

\begin{remglyph}
\hooglig{Center} the \remcom{large} \remcom{crown}.\\\hline
\end{remglyph}

\begin{prons}
\pronsrtwocone & \textbf{オウ (\={O})} & --- \\\hline
\pronsrthreecone & --- & --- \\\hline
\end{prons}

\begin{remprons}
At the \hooglig{center} of my being is an \comon{awesome} God.\\\hline
\end{remprons}

%{\middel}
% &  ()\\\hline
\begin{somme}
中央  & middle + center = \textit{center}  & チュウ{\middel}オウ \mbox{(CH\={U}{\middel}\={O})}  \\\hline
中央口  & middle + center + mouth = \textit{central exit/entrance} &  チュウ{\middel}オウ{\middel}ぐち \mbox{(CH\={U}{\middel}\={O}{\middel}guchi)}  \\\hline
\end{somme}
\end{multicols}

\begin{decomp}{5}
中央口&で&秋山さん&を&見ました。\\\hline
チュウオウぐち&で&あきやまさん&を&みました\\\hline
central entrance&で&Akiyama-san&を&saw\\\hline
\multicolumn{5}{|r|}{\textit{(He) saw Akiyama-san at the central entrance.}}\\\hline
\end{decomp}
%%--------------------------------------------------------------------
%%--------------------------------------------------------------------

\pagebreak